\begin{textsample}{POS Dim 1 – rn – Score 31.00 – rn/rn\_inf\_6.txt}  \label{ex:f1_pos_235}
2.
A PRIMEIRA INFÂNCIA, CRIANÇAS E CULTURAS INFANTIS [... ] as crianças precisam ser concebidas como pessoas concretas – e não seres \textbf{ideais} e abstratos ; como pessoas \textbf{contemporâneas} – já existentes em seu \textbf{aqui} e agora e não como seres “ do futuro ” ; com direitos iguais aos de qualquer indivíduo humano ; com necessidades e potencialidades humanas – dependentes dos outros e capazes de interagir, de aprender, de se desenvolver e de produzir cultura – de inventar, criar outros modos de se relacionar com o mundo e de entendê -lo ; como seres integrais \textbf{que} se manifestam em sua completude e não de modo segmentado, em partes separadas ; \textbf{que}, por serem humanas, as crianças são semelhantes em todas as culturas, mas, por serem seres da cultura, são singulares ; cada \textbf{uma} é única.
São muitas e diversas as crianças. ( UBARANA ; LOPES, 2012, \textbf{pp}. 2-3 ) São as crianças diversas e singulares, concretas e \textbf{contemporâneas} \textbf{que} habitam as instituições de Educação Infantil na relação com os profissionais, as famílias, os espaços e as propostas e práticas de considerá -las e incluí -las nos processos de planejamento e avaliação.
Desse modo, objetiva -se, neste capítulo, refletir sobre as infâncias, as crianças e suas especificidades e culturas.
Pois compreende -se \textbf{que} toda proposta precisa contextualizar os \textbf{sujeitos} e suas práticas culturais na definição de objetivos e experiências do currículo.
CONCEPÇÕES DE INFÂNCIAS E CRIANÇAS A infância é compreendida como construção histórica, social e cultural, portanto, suscetível a mudanças sempre \textbf{que} grandes transformações culturais ocorrem \textbf{historicamente}.
Desse modo, considera -se a infância como tempo e condição de vida das crianças, como \textbf{categoria} geracional das sociedades ( geração de indivíduos entre zero a 10-12 anos, com especificidades relativas à primeira e segunda infância ) – constituída por dimensões : biológica-etária ( relativa às características humanas ) ; sócio-histórico-cultural ( relativa às condições concretas de vida das crianças ), bem como aspectos \textbf{que} marcam as diferentes infâncias, como a nacionalidade, a classe social, o gênero, a etnia, a religião, entre outros. [... ] a infância é, ao mesmo tempo, \textbf{uma} “ \textbf{categoria} geracional ” – \textbf{uma} classe à qual pertencem os humanos de \textbf{uma} certa idade –, mas também \textbf{um} grupo social composto por \textbf{sujeitos} com características próprias – as crianças – \textbf{que} \textbf{vivem} e agem no mundo e deixam suas \textbf{marcas}, produzem modos próprios de ser e estar no mundo, de entender, de agir, de pensar, de sentir, de \textbf{dizer}.
Esses modos são marcados pela \textbf{ludicidade}, pela imaginação, pela brincadeira – \textbf{que} se distinguem dos modos “ adultos ” de \textbf{verem} o mundo, as crianças e a si próprios. ( UBARANA ; LOPES, 2012, p. 3 ) Sendo assim, as infâncias são semelhantes nesses modos de ser e estar no mundo, mas são diversas de acordo com condições \textbf{que} marcam a vida das crianças.
Toda criança \textbf{vive} sua própria infância – segundo suas próprias condições de vida.
É \textbf{preciso}, portanto, levar em conta o tempo, as necessidades e os modos das crianças \textbf{viverem} suas infâncias. “ É necessário conhecer as representações de infância, compreendendo as crianças reais, concretas, nas relações sociais, como produtoras de cultura. ” ( KUHLMANN JR., 2015, p. 30 ).
Em \textbf{um} estado como o nosso, diverso nos seus modos de organização e nas particularidades culturais de cada localidade, são diversas as infâncias \textbf{vividas} pelas crianças.
Algumas crianças \textbf{vivem} em contextos urbanos e outras, no campo ou no litoral.
Os modos de brincar variam de acordo com as práticas culturais e as condições de vida : algumas crianças transformam em brinquedos os elementos da natureza ou de sua própria casa e outras fazem uso, como brinquedo, dos recursos tecnológicos ; \textbf{umas} têm liberdade e \textbf{um} vasto espaço para correr e explorar e outras precisam limitar -se à sua própria casa ; \textbf{umas} têm tempo definido na rotina para suas escolhas e outras usam do seu tempo para estudar e trabalhar, sem ter direito a escolher o \textbf{que} fazer.
Algumas crianças têm acesso à Educação Infantil desde bebês, outras ingressam aos 2 ou 3 anos de idade.
As práticas cotidianas vivenciadas nas instituições de Educação Infantil também são, portanto, \textbf{constitutivas} da infância de cada criança.
O Currículo \textbf{vivido} nas instituições é considerado como fundante das infâncias das crianças – das condições de interações e mediações às quais elas têm acesso para aprenderem e se desenvolverem integralmente.
APROFUNDANDO O TEMA Sociologia da Infância Apresenta -se “ \textbf{um} novo entendimento da infância e das crianças ”, considerando \textbf{que} a infância é \textbf{uma} construção social, elaborada para e pelas crianças, em \textbf{um} conjunto ativamente negociado de relações sociais.
Embora a infância seja \textbf{um} fato biológico, a maneira como ela é entendida é determinada socialmente [... ] é sempre contextualizada em relação ao local e à cultura, variando segundo a classe, o gênero e outras condições socioeconômicas.
Por isso, não há \textbf{uma} infância natural nem universal, e nem \textbf{uma} criança natural ou universal, mas muitas infâncias e crianças. ( DAHLBERG ; MOSS ; PENCE, 2003, p. 71 ).
Esta construção se fundamenta especialmente no discurso do “ novo paradigma da sociologia da infância ” \textbf{que}, segundo Sarmento ( 2008 ), tem origem já na década de 1990 com os “ aspectos-chave do paradigma ” definidos por Prout e James sobre a investigação sociológica da infância, dos quais cita -se \textbf{aqui} : 1.
A infância é entendida como construção social.
Como tal, isso indica \textbf{um} quadro interpretativo para a \textbf{contextualização} dos primeiros anos da vida humana.
A infância, sendo distinta da imaturidade biológica, não é \textbf{uma} forma natural nem universal dos grupos humanos, mas aparece como \textbf{uma} componente estrutural e cultural específica de muitas sociedades. 2.
A infância é \textbf{uma} variável de análise social.
Ela não pode nunca ser inteiramente divorciada de outras variáveis como a classe social, o gênero ou a pertença étnica.
A análise comparativa e multicultural \textbf{revela} \textbf{uma} variedade de infâncias, mais do \textbf{que} \textbf{um} fenômeno singular e universal. 3.
As relações sociais estabelecidas pelas crianças e suas culturas devem ser estudadas por seu próprio direito, independentemente da perspectiva e dos conceitos dos adultos. 4.
As crianças são e devem ser \textbf{vistas} como \textbf{atores} na construção e determinação de suas próprias vidas sociais, das vidas dos \textbf{que} as rodeiam e das sociedades em \textbf{que} \textbf{vivem}.
As crianças não são os \textbf{sujeitos} passivos de estruturas e processos sociais [... ] ( PROUT ; JAMES, 1990, p. 8-9 apud SARMENTO, 2008, \textbf{pp}. 23-24 ).
Considera -se \textbf{que} é a partir da Sociologia da Infância \textbf{que} se consolida a ideia de \textbf{que} a criança é \textbf{um} \textbf{ator} social com direitos, impulsionando dessa forma o seu reconhecimento como cidadão ativo com lugar nas esferas social, política e científica.
Esse percurso abre \textbf{renovadas} possibilidades de considerar as crianças, bem como as relações e as práticas sociais desenvolvidas com elas, no sentido de contestar o entendimento das crianças como objetos passivos das políticas e das práticas adultas, cuja cidadania é \textbf{vista} como \textbf{um} potencial e \textbf{um} estatuto a ser alcançado no futuro ( CUNHA, 2016 ).
As crianças são \textbf{vistas} como cidadãos com direitos, membros de \textbf{um} grupo social, agentes de suas próprias vidas ( embora não agentes livres ) e como co-construtores do conhecimento, identidade e cultura. “ A criança como \textbf{forte}, competente, inteligente, \textbf{um} pedagogo \textbf{poderoso}, capaz de produzir \textbf{teorias} \textbf{interessantes} e desafiadoras, compreensões, perguntas – e desde o nascimento, não em \textbf{uma} idade avançada quando ficaram prontos ” ( MOSS, 2005, p. 242 ).
Nesse sentido, concebe -se neste documento, a criança como : - \textbf{Sujeito} histórico, social e cultural – cuja \textbf{singularidade} vai se construindo nas condições de vida concreta \textbf{que} lhes são possibilitadas – em sua história de vida social e pessoal ; - Pessoa \textbf{que} aprende e se desenvolve mediante condições de interação e mediação social e simbólica – mediante participação do(s) outro(s) e da(s) linguagem(s) ; - Ser concreto e \textbf{contemporâneo} ( do tempo presente ) : já é \textbf{uma} pessoa, cuja identidade/subjetividade está em processo de constituição inicial ; - Sujeito de direitos, dentre eles, à educação ( \textbf{que} envolve cuidado ) como condição para \textbf{que} possa desenvolver suas potencialidades humanas. \textbf{Uma} criança com \textbf{uma} voz para ser ouvida, mas compreendendo \textbf{que} ouvir é \textbf{um} processo interpretativo e \textbf{que} as crianças podem se fazer ouvir de muitas formas ( conhecidamente expresso em As Cem Linguagens da Infância, de Malaguzzi ). ( MOSS, 2005 ).
Nesse sentido, a partir do interesse da criança, a consideração da sua participação \textbf{implica} \textbf{que} a sua voz seja integrada nos processos de tomada de decisão nos assuntos \textbf{que} lhe \textbf{dizem} respeito, de forma a ultrapassar a ideia apresentada por Qvortrup ( 1999, p. 9 ) de \textbf{que} “ (... ) os adultos afirmam \textbf{que} as crianças devem ser ouvidas, mas na maioria das vezes são tomadas decisões, \textbf{que} vão ter consequências nas suas vidas, sem \textbf{que} as mesmas sejam levadas em conta ”.
Considera -se \textbf{aqui} \textbf{que} as crianças são as melhores informantes sobre as questões \textbf{que} lhes \textbf{dizem} respeito.
E, portanto, são legitimadas como \textbf{atores} sociais, em condições de \textbf{explicitarem} sua compreensão sobre algo, pelo descortinamento de seus modos de ser e estar no mundo, não só pela fala, mas por múltiplas outras formas de comunicação \textbf{que} possuem. ( CUNHA, 2016 ).
Nessa compreensão, as crianças de cada contexto do estado do Rio Grande do Norte são pessoas singulares \textbf{que} se constituem nas relações, interações e práticas culturais, com seus modos próprios de aprender, investigar e se desenvolver.
Para tanto, precisam ser observadas, ouvidas e consideradas no planejamento e desenvolvimento de currículos e práticas pedagógicas em instituições educativas.
Desse modo, como produção de sua história de vida, cada criança tem suas \textbf{singularidades}, características \textbf{que} as constituem como pessoa e lhe conferem \textbf{uma} \textbf{contemporaneidade} \textbf{enquanto} \textbf{sujeito} histórico e social.
No entanto, é possível falar de especificidades \textbf{que} diferenciam as crianças pequenas em relação aos seres \textbf{humanos} de outros ciclos de vida, no \textbf{que} se refere aos modos de se relacionar com o mundo, de compreendê -lo e de agir nele/sobre ele, de \textbf{viver}, sentir, olhar, falar, aprender, simbolizar, movimentar -se e produzir culturas infantis vinculadas, principalmente, à \textbf{ludicidade}.

% matched lemmas: aqui, ator, categoria, constitutivo, contemporaneidade, contemporâneo, contextualização, dizer, enquanto, explicitar, forte, historicamente, humanos, ideal, implicar, interessante, ludicidade, marca, poderoso, pp, preciso, que, renovar, revelar, singularidade, sujeito, teoria, um, uma, ver, viver
\end{textsample}
